\chapter{Ab initio ligand-field analysis (menu 38)}
\label{ch:menu38}
\menulabel{38}

\section{Purpose}

The AILFT module of ORCA (manual section 3.13.16) fits the parameters of a
ligand-field model -- a $5\times5$ or $7\times7$ one-electron matrix plus
the electron-repulsion parameters and the SOC constant -- to an ab initio
effective Hamiltonian built from SA-CASSCF (and NEVPT2) wave functions.
This menu reads the fitted parameters the engine already printed: the
ligand-field splitting, the Slater-Condon and Racah parameters per
theoretical level, the fit quality, and the SOC constant.  It never
refits.

\section{What you need}

An ORCA CASSCF output with the AILFT driver requested: \code{ActOrbs
dOrbs}/\code{fOrbs} (or the \code{LFTCase} shorthand, which also fixes the
principal quantum number), a CASSCF run over the $d^n$ or $f^n$ space, and
-- for the correlation-level parameters -- \code{!NEVPT2}.  Optionally
name a free-ion reference from the built-in table (the covered ions are
listed below) to get the nephelauxetic ratios level by level, or pass the
two values (free-ion Racah $B_0$ and SOC constant $\zeta_0$) by hand; a
hand-entered value overrides the table.

\section{Walkthrough}

The example reads the Ni$^{2+}$ $d^8$ fixture: the CASSCF-level parameters
(Racah $B$ = 1328.1~cm$^{-1}$, $C/B$ = 3.663), the NEVPT2 level
($B$ = 1218.0), the fit quality (CASSCF-level RMS 0.0~cm$^{-1}$ -- intrinsic
-- versus NEVPT2-level 457.4~cm$^{-1}$) and the SOC constant
ZETA\_D = 664.14~cm$^{-1}$.

\lstinputlisting[style=fbkcode, firstline=54, lastline=83,
  title={The menu-38 example (the two levels, the fit quality and the SOC
  constant; the built-in free-ion table supplies $B_0$ level by level, and
  the free ion against itself reads $\beta = 1.000$)}]
  {examples/expected/m38_ailft.txt}

\section{What you get}

The report \code{<output>.ailft.fbk.md}: the configuration header
($d^n$/$f^n$, CI blocks, MO range, centre atom); per level the
ligand-field eigenfunction energies and spread, the Slater-Condon
parameters (F0 marked \code{(fixed)} when from the raw two-electron
integrals; f shells carry F6), the Racah parameters and $C/B$, and the
\code{*.lft.gbw} file name; the fit quality (total and per-block RMS,
Pearson); the SOC constant; and -- when given -- the nephelauxetic ratios
$\beta = B/B_0$ and $\zeta/\zeta_0$.

\section{Conventions, cross-checks and boundaries (measured)}

\begin{readbox}
\begin{itemize}
  \item \textbf{The CASSCF-level RMS is intrinsic, not a quality gauge}:
    the LFT parametrization is exact for that level (measured: 0.0~cm$^{-1}$
    on the $d^8$ free ion); the correlation-level RMS reflects, among
    other things, the neglected anisotropy of electron-electron repulsion
    in covalent complexes (the rmsd discussion of Lang, Atanasov \& Neese
    2020).  Residual electron repulsion is fitted by $B$ and $C$ (or
    $F^2$, $F^4$, $F^6$); at the NEVPT2 level F0 (and Racah $A$) are not
    defined and are not printed (manual section 3.13.16).
  \item \textbf{Nephelauxetic readings resolve per level and per
    parameter}: the reductions relative to free ions are the classic
    covalency indicators (Jung, Atanasov \& Neese 2017, the
    actinide/lanthanide AILFT reference).  The built-in table covers the
    trivalent lanthanide and actinide series (Ce--Yb, Th--No; the
    published free-ion AILFT values of that paper, SI sections
    S3.5.1/S3.1.1: $F^2$/$F^4$/$F^6$ and $\zeta$ per parameter level),
    plus one measured entry (Ni$^{2+}$ $d^8$, frozen under
    \file{fixtures/ailft}).  The report resolves the ratio per parameter
    -- $F^2/F^2_0$ for f shells, $B/B_0$ for d shells,
    $\zeta/\zeta_0$ for both -- and a hand-entered value overrides the
    table.  The two kinds of entry were cross-checked against each other
    on Nd$^{3+}$: the NEVPT2-level $F^2$ and the $\zeta$ agree to better
    than 2\% between the published SARC2/DKH2 value and this toolkit's
    def2-SVP probe (the CASSCF-level $F^2$ carries a 7\% basis-set
    spread, recorded in the module; the Tb entry's NEVPT2 row repeats
    Gd's digits as printed in the source and is kept verbatim).  One
    measured engine boundary shaped the probes and is worth knowing: an
    ORCA 6.1.1 run that used \code{!TRAH} aborts inside the AILFT driver
    (\code{failed to retrieve the FAO matrix}), so the free-ion probes
    run without it.
  \item \textbf{The SOC constant comes from a CASSCF-orbital fit}
    (``SPIN ORBIT COUPLING (based on CASSCF orbitals)''); the confidence
    intervals of the fit are printed by the engine and are not restated.
  \item \textbf{The parameters are model quantities} of the ligand-field
    Hamiltonian -- this menu reads the engine's fit and never refits; the
    interpretation (AOM-style bonding analysis, two-shell problems) is
    the model's business, and the two-shell parameter spaces are covered
    by the manual's LFTCase list.
  \item \textbf{An unconverged run is flagged, not silently read}: when the
    output's own verdict is that the wavefunction is not fully converged
    (ORCA aborts such runs without using them for properties), the report
    opens with a RUN STATUS line marking the parameters diagnostic, and an
    energy-only convergence marker (the gradient criterion silent) gets the
    provisional caution -- the measured positives are the Dy(3+) free-ion
    probe's first attempt and a marker-rewritten probe variant.
  \item \textbf{Boundary (termination)}: LFDFT (ligand-field DFT) lives
    inside ADF, a commercial package; the AILFT
    route of this chapter covers the same analysis needs from open programs.
\end{itemize}
\end{readbox}
